\documentclass[preprint,12pt]{elsarticle}

\usepackage[T1]{fontenc}
\usepackage{lmodern}
\usepackage{microtype}
\usepackage{amsmath,amssymb,amsthm,mathtools}
\usepackage{booktabs,longtable,array,tabularx}
\usepackage{enumitem}
\usepackage{xcolor}
\usepackage[colorlinks=true,linkcolor=blue!55!black,citecolor=blue!55!black,urlcolor=blue!55!black]{hyperref}

\journal{Linear Algebra and its Applications}
\biboptions{sort&compress}

\newtheorem{theorem}{Theorem}[section]
\newtheorem{lemma}[theorem]{Lemma}
\newtheorem{proposition}[theorem]{Proposition}
\newtheorem{corollary}[theorem]{Corollary}
\theoremstyle{definition}
\newtheorem{definition}[theorem]{Definition}
\newtheorem{remark}[theorem]{Remark}

\newcommand{\CC}{\mathbb C}
\newcommand{\RR}{\mathbb R}
\newcommand{\NN}{\mathbb N}
\newcommand{\Sn}{\mathfrak S}
\newcommand{\per}{\operatorname{per}}
\newcommand{\tr}{\operatorname{tr}}
\newcommand{\diag}{\operatorname{diag}}
\newcommand{\Rea}{\operatorname{Re}}
\newcommand{\Ima}{\operatorname{Im}}
\newcommand{\adj}{\operatorname{adj}}
\newcommand{\rank}{\operatorname{rank}}
\newcommand{\ip}[2]{\left\langle #1,#2\right\rangle}
\newcommand{\norm}[1]{\left\lVert #1\right\rVert}
\newcommand{\abs}[1]{\left\lvert #1\right\rvert}
\newcommand{\mat}[1]{\begin{pmatrix}#1\end{pmatrix}}
\newcommand{\psd}{\succeq 0}
\newcommand{\pd}{\succ 0}
\newcommand{\eps}{\varepsilon}

\hypersetup{
  pdftitle={Permanental Dominance for Generalized Matrix Functions of Order Four},
  pdfauthor={Siwei Zeng},
  pdfsubject={Permanental dominance for all subgroup characters of order four}
}

\begin{document}

\begin{frontmatter}

\title{Permanental Dominance for Generalized Matrix Functions of Order Four}

\author[independent]{Siwei Zeng\corref{cor1}}
\ead{endlesslethe@gmail.com}
\ead[url]{https://orcid.org/0009-0008-8914-7143}
\cortext[cor1]{Corresponding author}
\affiliation[independent]{
  organization={Independent Researcher},
  orp={}
}

\begin{abstract}
For a subgroup $H\leq \Sn_4$ and an irreducible complex character $\chi$ of
$H$, let $d_\chi^H$ denote the associated generalized matrix function.  We
prove that $d_\chi^H(A)/\chi(1)\leq\per A$ for every $4\times4$ Hermitian
positive-semidefinite matrix $A$.  Thus permanental dominance holds in order
four for all subgroup characters, not only for the five irreducible
immanants associated with $\Sn_4$.  After reducing to correlation matrices,
we classify the thirty-seven irreducible-character cases arising from the
eleven conjugacy classes of subgroups of $\Sn_4$.  Thirty-five cases follow
from principal-minor, moment, and block-contraction inequalities.  The two
remaining conjugate characters of $A_4$ reduce to a polynomial inequality on
the cone of $3\times3$ positive-semidefinite Gram matrices.  Boundary
sum-of-squares, an exact interior Hessian calculation, and rational Gram
certificates establish its nonnegativity.
\end{abstract}

\begin{keyword}
permanent \sep generalized matrix function \sep immanant \sep
positive-semidefinite matrix \sep character theory \sep Gram matrix
\MSC[2020] 15A15 \sep 20C15 \sep 05E10
\end{keyword}

\end{frontmatter}

\section{Introduction}

The determinant and permanent are the two best-known members of the family
of generalized matrix functions.  Schur's classical theorem places the
determinant below every normalized irreducible immanant on the Hermitian
positive-semidefinite cone \cite{schur1918}.  In the opposite direction,
Lieb conjectured that the permanent lies above every normalized immanant
\cite{lieb1966}; the formulation became known as the permanental-dominance
conjecture through the work of Merris \cite{merris1987}.  In its subgroup
form, for an irreducible character $\chi$ of $H\leq\Sn_n$, the conjectured
inequality is
\[
  \frac{1}{\chi(1)}
  \sum_{\sigma\in H}\chi(\sigma)\prod_{i=1}^n A_{i,\sigma(i)}
  \leq \per A
  \qquad (A\psd).
\]
Taking $H=\Sn_n$ recovers the usual irreducible-immanant statement.

The conjecture has generated a substantial literature on dominance
orderings and sufficient analytic mechanisms.  Representative advances
include the ordering of hook immanants \cite{heyfron1988}, inequalities for
generalized matrix functions and row-appending methods
\cite{pate1991,pate1998,pate1999}, and the spectral approach of Soules
\cite{soules1994}.  Broader accounts and historical surveys are given by
Cheon and Wanless \cite{cheonwanless2005}, Zhang \cite{zhang2016}, and
Wanless \cite{wanless2022}.  More recent work relates immanant inequalities
to weight-space methods \cite{jingliu2025}, while the immanant form of
permanental dominance has acquired an operational interpretation through
generalized bosonic bunching probabilities \cite{gellerknill2026}.

Most of the literature concerns the immanant specialization $H=\Sn_n$.
That specialization is known in several broad families and, in particular,
through order thirteen; these results do not by themselves cover characters
of proper subgroups \cite{zhang2016,wanless2022}.  The distinction matters
already in order four: $\Sn_4$ has five irreducible characters, whereas all
subgroups of $\Sn_4$ contribute thirty-seven irreducible-character cases,
including non-real characters.  To the best of our knowledge, a uniform
proof covering every subgroup $H\leq\Sn_4$ and every irreducible character of
$H$ has not previously been recorded.

This paper proves precisely that all-subgroup statement for every
$4\times4$ Hermitian positive-semidefinite matrix.  The proof has three main
components.  First, diagonal scaling reduces the problem to correlation
matrices, and subgroup classification reduces it to a finite list of
character cases.  Second, general scalar and block inequalities settle all
but the two conjugate non-real characters of $A_4$.  Third, the remaining
character gap is converted into one polynomial inequality on the cone of
$3\times3$ Gram matrices and proved by exact boundary and interior
certificates.

The organizing idea is to write the exceptional gap as
\[
 c_0(u,v,w)+q^*K(u,v,w)q.
\]
The block positivity constraint says that $B(u,v,w)-qq^*$ is a Gram matrix.
After dividing out the nonzero coordinates of $q$, this residual becomes the
Gram matrix $R$ of three vectors, and the exceptional gap factors as
\[
 |q_1q_2q_3|^2\mathcal P(R).
\]
The problem is therefore moved from a character sum on four indices to the
nonnegativity of one explicit polynomial on the cone of $3\times3$ Gram
matrices.

Section~\ref{sec:statement-reductions} fixes the normalization and removes
the diagonal parameters.  Section~\ref{sec:finite-reduction} proves that the
finite list of cases is complete, and Section~\ref{sec:elementary-cases}
handles the cases covered by the general inequalities.
Sections~\ref{sec:a4-geometry} and~\ref{sec:Ppositive} contain the geometric
argument.  The proof is assembled in Section~\ref{sec:completion}; exact
coefficient certificates and the Lean verification statement are kept in
the appendices so that they do not interrupt the main line of reasoning.

\section{Statement and general reductions}\label{sec:statement-reductions}

Let $n\geq1$, let $H\leq\Sn_n$, and let $\chi:H\to\CC$ be a character.
For $A=(A_{ij})\in M_n(\CC)$ put
\begin{equation}\label{eq:def-dchi}
 d_\chi^H(A):=\sum_{\sigma\in H}\chi(\sigma)
   \prod_{i=1}^n A_{i,\sigma(i)},
 \qquad
 \per A:=\sum_{\sigma\in\Sn_n}\prod_{i=1}^n A_{i,\sigma(i)}.
\end{equation}
The normalization in the permanental-dominance conjecture is
$d_\chi^H/\chi(1)$.  Notice that no complex conjugate is placed on
$\chi(\sigma)$ in~\eqref{eq:def-dchi}; this convention agrees with all
formulae below.

\begin{lemma}[Reality]\label{lem:reality}
If $A=A^*$ and $\chi$ is the character of a finite-dimensional complex
representation of $H$, then $d_\chi^H(A)$ and $\per A$ are real.
\end{lemma}

\begin{proof}
Every finite-group representation may be made unitary, and hence
$\overline{\chi(\sigma)}=\chi(\sigma^{-1})$.  Hermitian symmetry gives
\[
 \overline{\prod_i A_{i,\sigma(i)}}
 =\prod_i A_{\sigma(i),i}
 =\prod_j A_{j,\sigma^{-1}(j)}.
\]
Conjugating the sum and replacing $\sigma$ by $\sigma^{-1}$ proves the first
claim.  The second is the special case $H=\Sn_n$, $\chi=1$.
\end{proof}

\begin{theorem}[Permanental dominance in order four]\label{thm:main}
Let $H\leq\Sn_4$, let $V$ be an irreducible finite-dimensional complex
representation of $H$, and let $\chi$ be its character.  For every
$A\in M_4(\CC)$ with $A\psd$,
\begin{equation}\label{eq:main}
  \frac{d_\chi^H(A)}{\dim V}\leq \per A.
\end{equation}
\end{theorem}

Theorem~\ref{thm:main} concerns order four only; it makes no assertion in
arbitrary order.

\subsection{Reduction to correlation matrices}

The normalization step is elementary but important because it removes all
diagonal parameters without an invertibility hypothesis.

\begin{lemma}[Diagonal congruence]\label{lem:scaling}
Let $s_1,\dots,s_n\in\RR$, and let
\[
 D=\diag(s_1,\dots,s_n).
\]
Then
\[
 d_\chi^H(DAD)=\left(\prod_i s_i\right)^2d_\chi^H(A),
 \qquad
 \per(DAD)=\left(\prod_i s_i\right)^2\per A.
\]
\end{lemma}

\begin{proof}
For every permutation $\sigma$,
\[
 \prod_i(DAD)_{i,\sigma(i)}
 =\prod_i(s_iA_{i,\sigma(i)}s_{\sigma(i)})
 =\left(\prod_i s_i\right)^2\prod_iA_{i,\sigma(i)}.
\]
Summation proves both identities.
\end{proof}

\begin{lemma}[Correlation reduction]\label{lem:correlation-reduction}
It is enough to prove Theorem~\ref{thm:main} for positive-semidefinite
matrices whose diagonal entries are all one.
\end{lemma}

\begin{proof}
Suppose first that $A_{ii}>0$ for all $i$ and put
$s_i=\sqrt{A_{ii}}$.  Since a Hermitian matrix has real diagonal, the matrix
\[
 C=\diag(s_i^{-1})A\diag(s_i^{-1})
\]
is positive semidefinite and has diagonal one.  Moreover,
\[
 A=\diag(s_i)C\diag(s_i).
\]
Thus Lemma~\ref{lem:scaling} transfers the desired inequality from $C$ to
$A$.

If $A_{ii}=0$ for some $i$, positivity of the $2\times2$ principal minors
implies $A_{ij}=A_{ji}=0$ for every $j$.  Each permutation monomial then
contains a zero entry in row $i$, so both sides of~\eqref{eq:main} vanish.
\end{proof}

We therefore write a general $4\times4$ correlation matrix as
\begin{equation}\label{eq:corr-six}
 A(a,b,c,d,e,f)=
 \mat{
 1&a&b&c\\
 \bar a&1&d&e\\
 \bar b&\bar d&1&f\\
 \bar c&\bar e&\bar f&1
 }\psd.
\end{equation}

\section{Finite character reduction}\label{sec:finite-reduction}

\subsection{Subgroups and characters}

Up to conjugacy, the subgroups of $\Sn_4$ have the following eleven types.
Here the superscripts ${\rm t},{\rm d},{\rm n}$ denote transposition,
double transposition, and normal, respectively; in the dihedral row,
$r=(1234)$ and $s=(13)$.
\begin{center}
\small
\begin{tabularx}{\textwidth}{@{}l>{\raggedright\arraybackslash}Xlrr@{}}
\toprule
type & representative & group & $|H|$ & $|\operatorname{Irr}(H)|$\\
\midrule
$1$ & $\{1\}$ & $1$ & 1&1\\
$C_2^{\rm t}$ & $\langle(12)\rangle$ & $C_2$ &2&2\\
$C_2^{\rm d}$ & $\langle(12)(34)\rangle$ & $C_2$ &2&2\\
$C_3$ & $\langle(123)\rangle$ & $C_3$ &3&3\\
$C_4$ & $\langle(1234)\rangle$ & $C_4$ &4&4\\
$V_4^{\rm n}$ & $\{1,(12)(34),(13)(24),(14)(23)\}$ & $C_2^2$ &4&4\\
$V_4^{\rm d}$ & $\langle(12),(34)\rangle$ & $C_2^2$ &4&4\\
$S_3$ & $\operatorname{Stab}(4)$ & $S_3$ &6&3\\
$D_8$ & $\langle r,s\rangle$ & $D_8$ &8&5\\
$A_4$ & $A_4$ & $A_4$ &12&4\\
$S_4$ & $S_4$ & $S_4$ &24&5\\
\bottomrule
\end{tabularx}
\end{center}
The final column sums to $37$.

\begin{lemma}[Subgroup classification]\label{lem:subgroup-classification}
Every subgroup of $\Sn_4$ is conjugate to exactly one of the eleven types in
the preceding table.
\end{lemma}

\begin{proof}
An intransitive subgroup preserves its orbit partition.  For partition
$2+1+1$ it is contained in $S_2$ and gives $1$ or $C_2^{\rm t}$.  For
$2+2$ it is contained in $S_2\times S_2$ and gives, up to conjugacy,
$1,C_2^{\rm t},C_2^{\rm d}$, or $V_4^{\rm d}$.  For $3+1$ it is contained
in the displayed point stabilizer $S_3$ and gives
$1,C_2^{\rm t},C_3$, or $S_3$.  This accounts for every intransitive type.

If the action is transitive, orbit--stabilizer shows that $4$ divides
$|H|$, while Lagrange's theorem gives $|H|\mid24$.  A transitive group of
order four is regular and is either $C_4$ or the normal
double-transposition group $V_4^{\rm n}$.  A group of order eight is a Sylow
$2$-subgroup, hence is conjugate to
$\langle(1234),(13)\rangle\cong D_8$.  The cases of order twelve and
twenty-four are $A_4$ and $S_4$, respectively.  The listed types are
distinguished by order, orbit structure, and, for the two order-two and two
Klein-four types, by cycle type.
\end{proof}

The character data used below are recorded in
Appendix~\ref{app:character-data}.  Orthogonality of the character-table rows
and the identity $\sum_\chi\chi(1)^2=|H|$ prove that no irreducible character
is missing.  For the two exceptional $A_4$ characters we fix
\begin{equation}\label{eq:omega}
 \omega=-\frac12-\frac{\sqrt3}{2}i,
 \qquad \omega^2+\omega+1=0.
\end{equation}

\subsection{Conjugation invariance}

\begin{lemma}[Transport]\label{lem:transport}
Let $g\in\Sn_4$, $K=g^{-1}Hg$, and let
$\chi^g(g^{-1}\sigma g)=\chi(\sigma)$.  If
$A^g_{ij}=A_{g(i),g(j)}$, then
\[
 d_{\chi^g}^{K}(A^g)=d_\chi^H(A),\qquad \per A^g=\per A,
 \qquad A\psd\Longrightarrow A^g\psd.
\]
\end{lemma}

\begin{proof}
Simultaneous permutation of rows and columns preserves positivity and the
permanent.  For a permutation monomial, changing $i$ to $g(i)$ gives
\[
 \prod_i A^g_{i,(g^{-1}\sigma g)(i)}
 =\prod_j A_{j,\sigma(j)}.
\]
Summing proves the assertion.
\end{proof}

Thus it suffices to prove dominance in the thirty-seven irreducible-character
cases associated with the eleven fixed representatives.

\section{Character cases from general inequalities}\label{sec:elementary-cases}

\subsection{Four cycle aggregates}

For~\eqref{eq:corr-six} define
\begin{align}
 \mathsf T={}&|a|^2+|b|^2+|c|^2+|d|^2+|e|^2+|f|^2,\label{eq:T}\\
 \mathsf D={}&|a|^2|f|^2+|b|^2|e|^2+|c|^2|d|^2,\label{eq:D}\\
 \mathsf C={}&2\Rea(a d\bar b+a e\bar c+b f\bar c+d f\bar e),\label{eq:C}\\
 \mathsf F={}&2\Rea(a d f\bar c+a e\bar f\bar b+b\bar d e\bar c).
 \label{eq:F}
\end{align}
Enumerating the cycle types of $\Sn_4$ gives
\begin{equation}\label{eq:per-aggregate}
 \per A=1+\mathsf T+\mathsf D+\mathsf C+\mathsf F.
\end{equation}
The five normalized $S_4$ generalized matrix functions, in the order of
the characters in~\eqref{eq:s4-table}, are
\begin{align}
 p_{[4]}&=1+\mathsf T+\mathsf D+\mathsf C+\mathsf F,\nonumber\\
 p_{[1^4]}&=1-\mathsf T+\mathsf D+\mathsf C-\mathsf F,\nonumber\\
 p_{[31]}&=1+\tfrac13\mathsf T-\tfrac13\mathsf D-\tfrac13\mathsf F,\label{eq:s4-rows}\\
 p_{[211]}&=1-\tfrac13\mathsf T-\tfrac13\mathsf D+\tfrac13\mathsf F,\nonumber\\
 p_{[22]}&=1+\mathsf D-\tfrac12\mathsf C.\nonumber
\end{align}

\begin{lemma}[Scalar correlation inequalities]\label{lem:scalar-toolbox}
If $A$ in~\eqref{eq:corr-six} is positive semidefinite, then
\begin{gather}
 0\leq\mathsf T,\quad 0\leq\mathsf D,\quad
 \mathsf D\leq\frac{\mathsf T^2}{4},\quad
 -2\mathsf D\leq\mathsf F\leq2\mathsf D\leq\mathsf T,
 \label{eq:scalar-basic}\\
 \frac{\mathsf T^2}{2}\leq\mathsf T+\frac32\mathsf C,
 \qquad 1+\mathsf D\leq\per A.\label{eq:moment-bound}
\end{gather}
In addition, if
\begin{equation}\label{eq:triangle-data}
 U=|a|^2+|b|^2+|d|^2,\qquad z=ad\bar b,
\end{equation}
then
\begin{equation}\label{eq:triangle-character-bounds}
 U+3\Rea z-\sqrt3\Ima z\geq0,
 \qquad U+3\Rea z+\sqrt3\Ima z\geq0.
\end{equation}
Every $3\times3$ principal permanent is at most $\per A$.
\end{lemma}

\begin{proof}
The inequalities $|a|^2,\ldots,|f|^2\leq1$ follow from the $2\times2$
principal minors.  The first four estimates in~\eqref{eq:scalar-basic} are
then AM--GM and
$2|\Rea(xy)|\leq|x|^2+|y|^2$, paired over opposite edges.

For the moment inequality, let $\lambda_1,\dots,\lambda_4\geq0$ be the
eigenvalues of $A$.  Direct trace expansion gives
\[
 \sum\lambda_i=4,\quad
 \sum\lambda_i^2=4+2\mathsf T,\quad
 \sum\lambda_i^3=4+6\mathsf T+3\mathsf C.
\]
The exact identity
\[
 \left(\sum_i\lambda_i\right)\left(\sum_i\lambda_i^3\right)
 -\left(\sum_i\lambda_i^2\right)^2
 =\sum_{i<j}\lambda_i\lambda_j(\lambda_i-\lambda_j)^2
\]
implies the first inequality in~\eqref{eq:moment-bound}.  It also gives
$\mathsf C\geq(\mathsf T^2-2\mathsf T)/3$.  To prove the second inequality,
subtract $1+\mathsf D$ from~\eqref{eq:per-aggregate}.  If
$\mathsf T\leq2$, then
\[
 \mathsf T+\mathsf C+\mathsf F
 \geq \mathsf T+\frac{\mathsf T^2-2\mathsf T}{3}-2\mathsf D
 \geq \frac{\mathsf T(2-\mathsf T)}6\geq0.
\]
If $\mathsf T\geq2$, the moment bound gives $\mathsf C\geq0$, while
$2\mathsf D\leq\mathsf T$ and $\mathsf F\geq-2\mathsf D$ give the same
conclusion.  In particular, $1\leq\per A$.

We next prove~\eqref{eq:triangle-character-bounds}.  Put
$u=|a|^2$, $v=|b|^2$, $w=|d|^2$, $S=u+v+w$, and write $z=x+iy$.
Positivity of the principal $3\times3$ block gives
\[
 0\leq u,v,w\leq1,\qquad x^2+y^2=uvw,
 \qquad 1-S+2x\geq0.
\]
Also $27uvw\leq S^3$ by AM--GM.  If $S\leq9/4$, Cauchy--Schwarz gives
\[
 (3x\pm\sqrt3y)^2\leq12(x^2+y^2)
 \leq\frac49S^3\leq S^2,
\]
which proves both inequalities.  If $S\geq9/4$, then
$x\geq(S-1)/2$ and $|y|\leq1$, so
\[
 S+3x\pm\sqrt3y
 \geq\frac{5S-3}{2}-\sqrt3>0.
\]

It remains to prove the principal-permanent contraction.  Choose Gram
vectors $x_1,\ldots,x_4$ for $A$, write
$x\odot y=x\otimes y+y\otimes x$, and set
\[
 y_1=x_2\odot x_3,\qquad y_2=x_1\odot x_3,
 \qquad y_3=x_1\odot x_2,qquad
 v=\bar c y_1+\bar e y_2+\bar f y_3.
\]
The identity
\[
 \ip{x\odot y}{z\odot w}
 =2\bigl(\ip{x}{z}\ip{y}{w}+\ip{x}{w}\ip{y}{z}\bigr)
\]
and direct collection of the cycle types give
\[
 \norm v^2=2\bigl(\per A-per A[\{1,2,3\}]\bigr).
\]
Hence the displayed principal permanent is at most $\per A$; simultaneous
relabeling proves the other three cases.
\end{proof}

We also need two symmetric $2+2$ contractions.  Put
\begin{align}
 R_+&:=1+|b|^2+|e|^2+\mathsf D
       +2\Rea(a d f\bar c),\label{eq:Rplus}\\
 R_\wedge&:=(1-|b|^2)(1-|e|^2)+|af-c\bar d|^2.\label{eq:Rwedge}
\end{align}

\begin{lemma}[Block contractions]\label{lem:block-contractions}
For a positive-semidefinite correlation matrix,
\begin{equation}\label{eq:block-bounds}
 R_+\leq\per A,\qquad R_\wedge\leq\per A,\qquad
 1+|a|^2|f|^2+|c|^2|d|^2\leq\per A.
\end{equation}
\end{lemma}

\begin{proof}
First note the exact identity
\begin{equation}\label{eq:Rplus-square}
 R_+=(1+|b|^2)(1+|e|^2)+|af+c\bar d|^2.
\end{equation}
Put $r=(1-|b|^2)^{1/2}$ and
\[
 y_0=a+\bar d b,\quad y_1=\bar d r,qquad
 g_0=c+fb,\quad g_1=fr.
\]
If $Y=|y_0|^2+|y_1|^2$, $G=|g_0|^2+|g_1|^2$, and
$h=\bar y_0g_0+\bar y_1g_1$, expanding the definitions gives the certificate
\begin{equation}\label{eq:Rplus-gap}
 \per A-R_+=Y+G+2\Rea(\bar e h).
\end{equation}
Cauchy--Schwarz and $|e|\leq1$ imply
$|\Rea(\bar e h)|^2\leq YG$.  Since
$Y+G\geq2\sqrt{YG}\geq2|\Rea(\bar e h)|$, the right-hand side
of~\eqref{eq:Rplus-gap} is nonnegative.  This proves the first contraction.

For the exterior contraction, take Gram vectors $x_1,\ldots,x_4$ for $A$.
Exterior Cauchy--Schwarz applied to $x_1\wedge x_3$ and
$x_2\wedge x_4$ gives
\[
 |af-c\bar d|^2\leq(1-|b|^2)(1-|e|^2),
\]
where the inner product of the two exterior tensors is $af-c\bar d$.

We need two consequences of the first contraction.  Relabeling its split
twice and using~\eqref{eq:Rplus-square} yields
$(1+|a|^2)(1+|f|^2)\leq\per A$ and
$(1+|c|^2)(1+|d|^2)\leq\per A$.  Put
\[
 x=|a|^2|f|^2,\qquad y=|c|^2|d|^2,\qquad P=\per A.
\]
If $x\leq y$, then $x\leq|c|^2$ and
$1+x+y\leq(1+|c|^2)(1+|d|^2)\leq P$; if $y\leq x$, use the other matching
bound.  This proves the third inequality in~\eqref{eq:block-bounds}.

Finally let $p=|b|^2$, $q=|e|^2$, and
$s=|af-c\bar d|$.  The exterior bound gives
$s^2\leq(1-p)(1-q)$, hence $0\leq s\leq1$.  The triangle inequality gives
$s\leq\sqrt{x}+\sqrt{y}$.  The two matching bounds and AM--GM imply
\[
 (1+\sqrt{x})^2\leq P,\qquad (1+\sqrt{y})^2\leq P.
\]
Consequently $(1+s/2)^2\leq P$.  Moreover
\[
 (1-p)(1-q)+s^2\leq(1+s/2)^2,
\]
because $p+q-pq\geq0$ and $s-3s^2/4\geq0$.  Therefore
$R_\wedge\leq P$, completing the proof.
\end{proof}

\subsection{Character cases covered by the general inequalities}

The following proposition records the finite analytic part of the proof.
Its table makes clear that the only cases not covered are the two non-real
$A_4$ characters.

\begin{proposition}[The thirty-five elementary character cases]\label{prop:elementary-cases}
For every representative subgroup in Lemma~\ref{lem:subgroup-classification},
all normalized generalized matrix functions other than those associated with
$\chi_\omega$ and
$\bar\chi_\omega$ of $A_4$ are bounded above by $\per A$.
\end{proposition}

\begin{proof}
The cyclic, Klein-four, and $S_3$ expansions are as follows.  In the $C_3$
and $S_3$ cases put $z=ad\bar b$.
\begin{center}
\small
\begin{tabularx}{\textwidth}{@{}l>{\raggedright\arraybackslash}X@{}}
\toprule
group & normalized generalized matrix functions\\
\midrule
$1$ & $1$\\
$C_2^{\rm t}$ & $1\pm|a|^2$\\
$C_2^{\rm d}$ & $1\pm|a|^2|f|^2$\\
$C_3$ & $1+2\Rea z$,\quad $1-\Rea z\pm\sqrt3\Ima z$\\
$V_4^{\rm n}$ & $1\pm r\pm s\pm t$ with the four character sign patterns,\\
& $r=|a|^2|f|^2$, $s=|b|^2|e|^2$, $t=|c|^2|d|^2$\\
$V_4^{\rm d}$ & $(1\pm|a|^2)(1\pm|f|^2)$\\
$S_3$ & $1+|a|^2+|b|^2+|d|^2+2\Rea z$,\\
& $1-|a|^2-|b|^2-|d|^2+2\Rea z$,\quad $1-\Rea z$\\
\bottomrule
\end{tabularx}
\end{center}
These expressions follow by listing the elements of the corresponding
subgroup.  We give the inequalities explicitly.  Set
\[
 U=|a|^2+|b|^2+|d|^2,\quad z=x+iy,\quad
 P_3=1+U+2x.
\]
Lemma~\ref{lem:scalar-toolbox} gives $P_3\leq\per A$ and
$U+3x\pm\sqrt3y\geq0$.  Thus
\[
 P_3-(1-x\pm\sqrt3y)=U+3x\mp\sqrt3y\geq0,
\]
which verifies the $C_3$ cases; the value $1+2x$ is at most $P_3$ because
$U\geq0$.  For $S_3$, the trivial case is $P_3$, the sign case differs from
$P_3$ by $-2U$, and the standard case satisfies
$P_3-(1-x)=U+3x\geq0$, obtained by averaging the two triangle bounds.

The trivial and $C_2^{\rm t}$ cases are bounded by $P_3$ (or by $1$ after
choosing the negative sign).  For $C_2^{\rm d}$, the positive value is at
most $1+|a|^2|f|^2+|c|^2|d|^2\leq\per A$, while the negative value is at
most one.  For $V_4^{\rm n}$, the all-positive value is
$1+\mathsf D\leq\per A$ and each other sign pattern is no larger.  For
$V_4^{\rm d}$, all four values are no larger than
$(1+|a|^2)(1+|f|^2)\leq\per A$.

For $C_4$, put
\[
 r=|b|^2|e|^2,\qquad z=adf\bar c.
\]
Its four normalized functions are
\begin{equation}\label{eq:c4-rows}
 1+r+2\Rea z,\quad 1-r-2\Ima z,\quad
 1+r-2\Rea z,\quad 1-r+2\Ima z.
\end{equation}
Put $x=|a|^2|f|^2$ and $y=|c|^2|d|^2$.  The estimates
\[
 2|\Rea z|,\ 2|\Ima z|
 \leq x+y
\]
show that the two values containing $+r$ are at most
$1+r+x+y=1+\mathsf D\leq\per A$.  The other two are at most
$1-r+x+y\leq1+x+y\leq\per A$ by~\eqref{eq:block-bounds}.

For $D_8$, write $p=|b|^2$, $q=|e|^2$.  The five normalized functions are
\begin{align}
 R_+,&\quad
 (1+p)(1+q)-|af+c\bar d|^2,\nonumber\\
 &(1-p)(1-q)-|af-c\bar d|^2,\quad
 R_\wedge,\quad 1-pq.\label{eq:d8-rows}
\end{align}
The first and fourth are~\eqref{eq:block-bounds}.  By~\eqref{eq:Rplus-square},
the second differs from $R_+$ only by replacing
$+|af+c\bar d|^2$ with $-|af+c\bar d|^2$, and hence is no larger.  The third
and fifth are at most $1$, hence at most $\per A$.

For $S_4$, the sign case follows from
$\mathsf T+\mathsf F\geq\mathsf T-2\mathsf D\geq0$.  For the two remaining
nontrivial estimates, direct subtraction gives
\begin{align*}
 3(\per A-p_{[31]})
   &=2\mathsf T+4\mathsf D+3\mathsf C+4\mathsf F
     \geq\mathsf T^2-4\mathsf D\geq0,\\
 2(\per A-p_{[22]})
   &=2\mathsf T+3\mathsf C+2\mathsf F
     \geq\mathsf T^2-4\mathsf D\geq0.
\end{align*}
Here we used~\eqref{eq:scalar-basic}--\eqref{eq:moment-bound}.  Finally,
$p_{[211]}\leq1\leq\per A$ because
$\mathsf F\leq\mathsf T$ and $\mathsf D\geq0$.

Finally, the trivial $A_4$ case is the average of the trivial and sign $S_4$
cases, while its three-dimensional case is the average of the $[31]$ and
$[211]$ cases.  This proves the proposition.
\end{proof}

\section{The exceptional \texorpdfstring{$A_4$}{A4} geometry}
\label{sec:a4-geometry}

\subsection{The exceptional character gap}

It remains to treat the two conjugate non-real characters of $A_4$.  Their
common geometric content is isolated below.

Let
\begin{equation}\label{eq:Buvw}
 B(u,v,w)=\mat{1&u&v\\\bar u&1&w\\\bar v&\bar w&1}.
\end{equation}
Define
\begin{equation}\label{eq:c0}
 c_0(u,v,w)=|u|^2+|v|^2+|w|^2
 +2\Rea\bigl((1-\omega)uw\bar v\bigr)
\end{equation}
and the Hermitian matrix
\begin{equation}\label{eq:K}
 K(u,v,w)=
 \mat{
 1&(1-\omega^2)u+v\bar w&uw+(1-\omega)v\\
 (1-\omega)\bar u+w\bar v&1&(1-\omega^2)w+v\bar u\\
 \bar u\bar w+(1-\omega^2)\bar v&(1-\omega)\bar w+u\bar v&1
 }.
\end{equation}

\begin{lemma}[Character expansion]\label{lem:a4-expansion}
For the correlation matrix~\eqref{eq:corr-six},
\begin{equation}\label{eq:a4-gap}
 \per A-d_{\chi_\omega}^{A_4}(A)
 =c_0(a,b,d)+\mat{\bar c&\bar e&\bar f}
 K(a,b,d)\mat{c\\e\\f}.
\end{equation}
For $\bar\chi_\omega$, the same formula holds after an odd simultaneous
relabeling of the four indices.
\end{lemma}

\begin{proof}
Choose the labeling of $C_+$ so that $\chi_\omega(C_+)=\omega$.  Listing the
twelve elements gives the independently checkable expansion
\begin{align*}
d_{\chi_\omega}^{A_4}(A)
={}&1+|a|^2|f|^2+|b|^2|e|^2+|c|^2|d|^2\\
&+\omega\bigl(e\bar d\bar f+ad\bar b+bf\bar c+c\bar a\bar e\bigr)\\
&+\omega^2\bigl(df\bar e+ae\bar c+b\bar a\bar d+c\bar b\bar f\bigr).
\end{align*}
The four monomials on the last line are the conjugates of those on the
preceding line, in the inverse order, so the displayed quantity is real.
Subtracting it from~\eqref{eq:per-aggregate}, using
$\omega^2=-1-\omega$, and collecting the constant, diagonal, and off-diagonal
terms in $(c,e,f)$ yields~\eqref{eq:a4-gap}.  An odd permutation interchanges
the two three-cycle classes and therefore exchanges $\chi_\omega$ with its
conjugate.
\end{proof}

The remaining task is consequently the following four-vector inequality.

\begin{theorem}[Four-vector inequality]\label{thm:four-vector}
Suppose
\begin{equation}\label{eq:completion}
 \mathcal A=\mat{B(u,v,w)&q\\q^*&1}\psd,\qquad q\in\CC^3.
\end{equation}
Then
\begin{equation}\label{eq:four-vector}
 c_0(u,v,w)+q^*K(u,v,w)q\geq0.
\end{equation}
\end{theorem}

The proof is developed in the residual normalization below and in
Section~\ref{sec:Ppositive}.

\subsection{Residual-vector normalization}\label{sec:normalization}

There is a useful geometric meaning behind the next change of variables.
Realize the block matrix~\eqref{eq:completion} as the Gram matrix of four unit
vectors.  The vector corresponding to the last row and column determines a
common one-dimensional direction.  Subtracting the projections of the first
three vectors onto that direction leaves three residual vectors, whose Gram
matrix is $B-qq^*$.  Dividing the residual vectors by the coordinates of the
removed projections does not change positivity; it merely produces the
scale-free matrix $R$ below.  This is the step that turns a four-vector
completion problem into a problem on the cone of three-vector Gram matrices.

Assume first that $q_1q_2q_3\neq0$.  The Schur complement of the final
diagonal entry in~\eqref{eq:completion} is
$B-qq^*\psd$.  Congruence by
$\diag(\bar q_1^{-1},\bar q_2^{-1},\bar q_3^{-1})$ gives the residual Gram
matrix
\begin{equation}\label{eq:residual-gram}
 R=\mat{
 n_1&h_{12}&h_{13}\\
 \bar h_{12}&n_2&h_{23}\\
 \bar h_{13}&\bar h_{23}&n_3
 }\psd,
\end{equation}
where
\begin{equation}\label{eq:residual-coordinates}
 n_i=\frac{1-|q_i|^2}{|q_i|^2},\qquad
 h_{12}=\frac{u}{q_1\bar q_2}-1,\quad
 h_{13}=\frac{v}{q_1\bar q_3}-1,\quad
 h_{23}=\frac{w}{q_2\bar q_3}-1.
\end{equation}
Thus $R$ is the Gram matrix of three residual vectors $w_1,w_2,w_3$ with
$n_i=\norm{w_i}^2$ and $h_{ij}=\ip{w_i}{w_j}$.

We now define the normalized geometric polynomial.  Put
\begin{align}
 E={}&n_1n_2+|h_{12}|^2+n_1n_3+|h_{13}|^2+n_2n_3+|h_{23}|^2,\nonumber\\
 x_{12,13}={}&n_1h_{23}+h_{13}\bar h_{12},\nonumber\\
 x_{12,23}={}&h_{12}h_{23}+n_2h_{13},\nonumber\\
 x_{13,23}={}&n_3h_{12}+h_{13}\bar h_{23},\nonumber\\
 S={}&E+2\Rea(x_{12,13}+x_{12,23}+x_{13,23}),\label{eq:tensor-terms}\\
 S_\omega={}&E+2\Rea(\omega^2x_{12,13}+\omega x_{12,23}
                     +\omega^2x_{13,23}),\nonumber\\
 \tau={}&h_{13}\bar h_{12}\bar h_{23},\nonumber\\
 \Gamma={}&n_3|h_{12}|^2+n_2|h_{13}|^2+n_1|h_{23}|^2
              +3\Rea\tau+\sqrt3\Ima\tau.\nonumber
\end{align}
Finally set
\begin{equation}\label{eq:geometric-P}
 \mathcal P(R)=24+6\left(n_1+n_2+n_3
 +2\Rea(h_{12}+h_{13}+h_{23})\right)+2S-S_\omega+\Gamma.
\end{equation}

\begin{lemma}[Exact normalization identity]\label{lem:normalization-identity}
With~\eqref{eq:residual-coordinates},
\begin{equation}\label{eq:normalization-identity}
 c_0(u,v,w)+q^*K(u,v,w)q
 =|q_1q_2q_3|^2\mathcal P(R).
\end{equation}
\end{lemma}

\begin{proof}
Equations~\eqref{eq:residual-coordinates} give
\[
 u=q_1\bar q_2(1+h_{12}),\quad
 v=q_1\bar q_3(1+h_{13}),\quad
 w=q_2\bar q_3(1+h_{23}),
\]
and $|q_i|^2(1+n_i)=1$.  Let $Q=|q_1q_2q_3|^2$.  Dividing first by $Q$
makes every cancellation visible.  Direct substitution in~\eqref{eq:c0}
gives
\begin{align}
\frac{c_0(u,v,w)}{Q}={}&
 (1+n_3)|1+h_{12}|^2+(1+n_2)|1+h_{13}|^2
 +(1+n_1)|1+h_{23}|^2\nonumber\\
&+2\Rea\bigl((1-\omega)(1+h_{12})(1+h_{23})
                   (1+\bar h_{13})\bigr).\label{eq:c0-normalized}
\end{align}
For the quadratic part, the diagonal terms and the three upper-triangular
entries of $K$ give
\begin{align}
\frac{q^*Kq}{Q}={}&
 (1+n_2)(1+n_3)+(1+n_1)(1+n_3)+(1+n_1)(1+n_2)\nonumber\\
&+2\Rea\Bigl(
 (1-\omega^2)(1+n_3)(1+h_{12})
 +(1+h_{13})(1+\bar h_{23})\nonumber\\
&\qquad +(1+h_{12})(1+h_{23})
 +(1-\omega)(1+n_2)(1+h_{13})\nonumber\\
&\qquad +(1-\omega^2)(1+n_1)(1+h_{23})
 +(1+h_{13})(1+\bar h_{12})\Bigr).
\label{eq:K-normalized}
\end{align}
Equations~\eqref{eq:c0-normalized}--\eqref{eq:K-normalized} are obtained
term by term; for example,
\[
 \frac{|u|^2}{Q}=(1+n_3)|1+h_{12}|^2,qquad
 \frac{uw\bar v}{Q}=(1+h_{12})(1+h_{23})(1+\bar h_{13}).
\]
Expanding these two displayed identities and using
$1+\omega+\omega^2=0$ groups the degree-two terms as $2S-S_\omega$ and the
degree-three terms as $\Gamma$; the remaining terms are the constant and
linear part in~\eqref{eq:geometric-P}.  Thus their sum is
$\mathcal P(R)$, which proves~\eqref{eq:normalization-identity}.
\end{proof}

Thus Theorem~\ref{thm:four-vector} follows from the next result.

\begin{theorem}[Geometric polynomial]\label{thm:Ppositive}
For every Hermitian matrix $R$ of the form~\eqref{eq:residual-gram},
\[
 R\psd\quad\Longrightarrow\quad\mathcal P(R)\geq0.
\]
\end{theorem}

\section{Positivity of the geometric polynomial}\label{sec:Ppositive}

The proof uses the geometry of the Gram cone in three stages.  On the
rank-one locus the three residual vectors are collinear and the polynomial
has a direct sum-of-squares form.  In the positive-definite region, two
residual vectors are used as coordinates and the third is minimized through
an explicit two-dimensional Hermitian Hessian.  Finally, every singular Gram
matrix is obtained as a limit of positive-definite ones.  The rank-one
calculation is also the algebraic base coefficient in the interior minimum
certificate, so it is both explanatory and logically active.

\subsection{The rank-one boundary}

Suppose first that $R$ has rank one.  Write
\[
 n_1=|x|^2,\quad n_2=|y|^2,\quad n_3=|z|^2,\qquad
 h_{12}=\bar xy,\quad h_{13}=\bar xz,\quad h_{23}=\bar yz.
\]
Define
\begin{align}
 R_0(x,y,z)={}&4+|x+y+z|^2+|xy+xz+yz|^2\nonumber\\
 &\quad-(|xy|^2+|xz|^2+|yz|^2)+|xyz|^2,\label{eq:R0}\\
 \Phi(x,y,z)={}&xy+\omega xz+\omega^2yz.
\end{align}

\begin{lemma}[Rank-one sum of squares]\label{lem:rank-one-sos}
For all $x,y,z\in\CC$,
\begin{equation}\label{eq:P-rank-one}
 \mathcal P(R)=2\bigl(3R_0(x,y,z)+|\Phi(x,y,z)|^2\bigr)\geq0.
\end{equation}
\end{lemma}

\begin{proof}
Regard $R_0$ as a quadratic in $z$.  Put
\[
 K_0=|1+\bar xy|^2,\quad
 L_0=\overline{x+y}+\overline{xy}(x+y),\quad
 C_0=4+|x+y|^2.
\]
Direct expansion gives
\[
 R_0=K_0|z|^2+2\Rea(L_0z)+C_0
\]
and the exact discriminant identity
\[
 C_0K_0-|L_0|^2=4(1+\Rea(\bar xy))^2\geq0.
\]
Completing the square proves $R_0\geq0$.  A second direct expansion using
$1+\omega+\omega^2=0$ gives the equality in~\eqref{eq:P-rank-one}; its
right-hand side is nonnegative.
\end{proof}

\subsection{The positive-definite region}

It remains to treat a positive-definite residual Gram matrix.  Write
\begin{equation}\label{eq:G-data}
 a=n_1>0,\quad b=n_2,\quad c=n_3,\quad
 h_{12}=g=r-is,\quad p=h_{13},\quad q=h_{23},
\end{equation}
and put
\begin{equation}\label{eq:Gram-G-delta}
 G=\mat{a&g\\\bar g&b},\qquad
 \delta=\det G=ab-r^2-s^2>0,\qquad h=\mat{p\\q}.
\end{equation}
There is a unique $\zeta\in\CC^2$ such that $h=G\zeta$.  Positivity of
$R$ tested on $(-\zeta,1)$ gives the Schur slack
\begin{equation}\label{eq:schur-slack}
 \zeta^*G\zeta\leq c.
\end{equation}

Set
\begin{align}
 \lambda={}&6+a+b+r^2+s^2+5r-\sqrt3s,\label{eq:lambda}\\
 P_0={}&24+6(a+b+2r)+ab+r^2+s^2,\label{eq:P0}\\
 \ell={}&\mat{
 6+(2-\omega)g+(2-\omega^2)b\\
 6+(2-\omega)a+(2-\omega^2)\bar g},\label{eq:ell}\\
 M={}&\mat{
 1+b&2-\omega+(1-\omega)g\\
 2-\omega^2+(1-\omega^2)\bar g&1+a},\label{eq:M}\\
 T={}&\lambda I_2+MG.\label{eq:Tmatrix}
\end{align}

\begin{lemma}[Quadratic decomposition]\label{lem:quadratic-decomposition}
With the preceding notation,
\begin{align}
 \mathcal P(R)
 ={}&P_0+2\Rea(\ell^*h)+\lambda c+\Rea(h^*Mh)\label{eq:P-quadratic}\\
 ={}&Q(\zeta)+\lambda(c-\zeta^*G\zeta),\nonumber
\end{align}
where
\begin{equation}\label{eq:Qzeta}
 Q(\zeta)=P_0+2\Rea((G\ell)^*\zeta)+\zeta^*(GT)\zeta.
\end{equation}
Moreover $GT$ is Hermitian.
\end{lemma}

\begin{proof}
Substitution of~\eqref{eq:G-data} in~\eqref{eq:geometric-P} gives the first
identity.  Replacing $h$ by $G\zeta$, using $G=G^*$ and
$T=\lambda I+MG$, gives the second.  Finally
$GT=\lambda G+GMG$ is Hermitian because $G$ and $M$ are Hermitian.
\end{proof}

\begin{lemma}[Positivity of the scalar part]\label{lem:lambda-positive}
Under $a,b\geq0$ and $r^2+s^2\leq ab$, one has $\lambda>0$.
\end{lemma}

\begin{proof}
Let $\rho=\sqrt{r^2+s^2}$.  Cauchy--Schwarz and AM--GM give
\[
 5r-\sqrt3s\geq-2\sqrt7\rho,\qquad a+b\geq2\rho.
\]
Consequently
\[
 \lambda\geq(\rho-(\sqrt7-1))^2+2(\sqrt7-1)>0.
\]
\end{proof}

The remaining algebra is organized by the Gram defect $\delta$.  Define
\begin{equation}\label{eq:Ndef}
 N=\Rea\left(P_0\det T-\ell^*G\adj(T)\ell\right).
\end{equation}

\begin{lemma}[Hessian and minimum certificates]\label{lem:hessian-certs}
For $a>0$ and $\delta\geq0$,
\begin{equation}\label{eq:detT-expansion}
 a^3\Rea\det T=D_0+\delta D_1+\delta^2D_2,
\end{equation}
where
\begin{align}
 D_0&=a^3\lambda_0H_0>0,\nonumber\\
 D_1&=a\bigl((4+11a)(r^2+s^2)+a(20+a)r\nonumber\\
 &\qquad\quad-\sqrt3a(4+3a)s
          +a(18+10a+3a^2)\bigr)>0,\label{eq:Dcoeffs}\\
 D_2&=a(a+1)(a+2)>0.
\end{align}
Here $\lambda_0$ denotes~\eqref{eq:lambda} with
$b=(r^2+s^2)/a$, and
\[
 H_0=2\lambda_0-6+4(r^2+s^2)>0.
\]
Furthermore
\begin{equation}\label{eq:N-expansion}
 a^3N=N_0+\delta N_1+\delta^2N_2+\delta^3N_3,
\end{equation}
with $N_0\geq0$ and $N_1,N_2,N_3>0$.
\end{lemma}

\begin{proof}
The determinant identity follows by expanding~\eqref{eq:Tmatrix} after
$b=(r^2+s^2+\delta)/a$.  To make all sign checks explicit, put
$\rho=(r^2+s^2)^{1/2}$ and $k=\sqrt7-1$.  On the boundary $\delta=0$, the
two estimates in Lemma~\ref{lem:lambda-positive} give
\[
 \lambda_0\geq(\rho-k)^2+2k.
\]
Consequently
\begin{equation}\label{eq:H0-lower}
 H_0\geq6\rho^2-4k\rho+6
 =6\left(\rho-\frac{k}{3}\right)^2+\frac{2+4\sqrt7}{3}>0.
\end{equation}
This proves $D_0>0$.

For $D_1$, set
\[
 Q=4+11a,\quad L_r=a(20+a),\quad
 L_s=-\sqrt3a(4+3a),\quad C=a(18+10a+3a^2).
\]
The identity
\begin{equation}\label{eq:two-variable-completion}
 4Q\bigl(Q(r^2+s^2)+L_rr+L_ss+C\bigr)
 =(2Qr+L_r)^2+(2Qs+L_s)^2+4QC-L_r^2-L_s^2
\end{equation}
reduces positivity to the exact residual
\[
 4QC-L_r^2-L_s^2
 =4a(72+126a+94a^2+26a^3)>0.
\]
Since $a,Q>0$, this proves $D_1>0$; the assertion for $D_2$ is immediate.

The expansion~\eqref{eq:N-expansion} is obtained in the same way from
\eqref{eq:Ndef}.  On the boundary $\delta=0$, the first two residual vectors
are collinear.  Define
\[
 x_0=\sqrt a,\qquad y_0=\frac{g}{\sqrt a},\qquad
 L_0'=x_0\ell_1+y_0\ell_2,
\]
where $\ell_1,\ell_2$ are the components of~\eqref{eq:ell}, evaluated at
$b=(r^2+s^2)/a$.  Then $|x_0|^2=a$, $|y_0|^2=b$, and
$\bar x_0y_0=g$.  Lemma~\ref{lem:rank-one-sos}, viewed as a quadratic in the
third scalar coordinate, gives
$|L_0'|^2\leq P_0H_0$.  Direct substitution in~\eqref{eq:Ndef} now yields
the exact boundary coefficient
\begin{equation}\label{eq:N0-certificate}
 N_0=a^3\lambda_0(P_0H_0-|L_0'|^2)\geq0.
\end{equation}

The remaining coefficients are explicit.  First,
\[
 N_3=(a+2)(a^2+6)>0.
\]
Next, collecting the terms in $(r,s)$ gives
\begin{align*}
 N_2={}&36(r^2+s^2)
 +a(120+84r-36\sqrt3s+36(r^2+s^2))\\
 &+a^2(60+24r-24\sqrt3s+16(r^2+s^2))\\
 &+a^3(16+2(r^2+s^2))+2a^4.
\end{align*}
Thus $N_2=Q_2(r^2+s^2)+L_{2r}r+L_{2s}s+C_2$, where
\begin{align*}
 Q_2&=36+36a+16a^2+2a^3,\\
 L_{2r}&=84a+24a^2,\qquad
 L_{2s}=-\sqrt3(36a+24a^2),\\
 C_2&=120a+60a^2+16a^3+2a^4.
\end{align*}
Applying~\eqref{eq:two-variable-completion} gives the strictly positive
residual
\[
 4Q_2C_2-L_{2r}^2-L_{2s}^2
 =16a(a^6+16a^5+112a^4+318a^3+588a^2+936a+1080)>0.
\]
Therefore $N_2>0$.

For $N_1$, put $X=r$, $Y=\sqrt3s$ and
$\nu=X^2+Y^2/3$.  Then
\begin{equation}\label{eq:N1-B}
 N_1=36\nu^2+aB_1+a^2B_2+a^3B_3+a^4B_4+6a^5,
\end{equation}
where
\begin{align*}
 B_1&=\nu(240+168X-72Y+54\nu),\\
 B_2&=432+624X-240Y+492X^2+108Y^2-168XY\\
    &\qquad +(156X-84Y)\nu+38\nu^2,\\
 B_3&=168+72X-168Y+54X^2+46Y^2-84XY\\
    &\qquad +(20X-24Y)\nu+6\nu^2,\\
 B_4&=60+24X-24Y+14\nu.
\end{align*}
Appendix~\ref{app:sos} gives exact certificates for
$B_1\geq0$, $B_2>0$, $B_4>0$, and
$B_2+aB_3+a^2B_4>0$.  Hence~\eqref{eq:N1-B} is strictly positive.
\end{proof}

\begin{proposition}[Positive-definite residual Gram matrices]
If $R\pd$, then $\mathcal P(R)>0$.
\end{proposition}

\begin{proof}
Equations~\eqref{eq:detT-expansion} and~\eqref{eq:Dcoeffs} show
$\Rea\det T>0$.  Direct multiplication gives the exact leading-entry
identity
\begin{equation}\label{eq:GT-leading}
 (GT)_{11}=aH_0+\delta(a+1)>0.
\end{equation}
Moreover
\[
 \det(GT)=\det G\det T=\delta\det T.
\]
Because $GT$ is Hermitian, its determinant is real; since $\delta>0$, this
also proves $\det T=\Rea\det T$.  Hence
$\det(GT)=\delta\Rea\det T>0$.  The two leading principal minors are
positive, so the $2\times2$ Hermitian matrix $GT$ is positive definite.

The critical-point equation for~\eqref{eq:Qzeta} is
$GT\zeta=-G\ell$.  Since $G$ is invertible, it is equivalent to
$T\zeta=-\ell$, and therefore
\[
 \zeta_0=-\frac{\adj(T)\ell}{\Rea\det T},
\]
and completing the Hermitian square gives
\[
 \min_\zeta Q(\zeta)=\frac{N}{\Rea\det T}.
\]
By~\eqref{eq:N-expansion}, $N>0$ when $\delta>0$.  Therefore
$Q(\zeta)>0$ for all $\zeta$.  Lemma~\ref{lem:lambda-positive} and the Schur
slack~\eqref{eq:schur-slack} now give
\[
 \mathcal P(R)=Q(\zeta)+\lambda(c-\zeta^*G\zeta)>0.
\]
\end{proof}

\subsection{Closure of the Gram cone}

\begin{proof}[Proof of Theorem~\ref{thm:Ppositive}]
For $R\psd$ and $t>0$, the matrix $R+tI$ is positive definite, and hence
$\mathcal P(R+tI)>0$ by the preceding proposition.  Since
$\mathcal P$ is a polynomial in the real and imaginary parts of the entries,
letting $t\downarrow0$ gives $\mathcal P(R)\geq0$.
\end{proof}

\begin{proof}[Proof of Theorem~\ref{thm:four-vector}]
If all coordinates of $q$ are nonzero, the residual matrix
\eqref{eq:residual-gram} is positive semidefinite.  Theorem~\ref{thm:Ppositive}
and~\eqref{eq:normalization-identity} give~\eqref{eq:four-vector}.

For arbitrary $q$, choose phases $p_i$ of modulus one, taking
$p_i=q_i/|q_i|$ when $q_i\neq0$ and $p_i=1$ otherwise.  Convexly interpolate
the block matrix~\eqref{eq:completion} with the rank-one correlation
completion generated by $(p_1,p_2,p_3,1)$.  Its new last column is
$q_i(t)=(1-t)q_i+tp_i$, which is nonzero for $0<t<1$.  The interpolated
matrix remains positive semidefinite, so the already proved case applies.
Letting $t\downarrow0$ and using continuity proves the result.
\end{proof}

\subsection{The resulting matrix certificate}

The four-vector inequality has a useful matrix consequence.  We record it
after the direct proof so that it does not obscure the residual-vector
argument that establishes it.

\begin{lemma}[Ellipsoid-to-adjugate certificate]\label{lem:ellipsoid}
Let $B,K\in M_m(\CC)$ with $B\psd$ and $K=K^*$, and let $c\in\RR$.  If
\[
 \mat{B&q\\q^*&1}\psd\quad\Longrightarrow\quad c+q^*Kq\geq0
\]
for every $q\in\CC^m$, then
\begin{equation}\label{eq:adj-certificate}
 (\det B)K+c\,\adj(B)\psd.
\end{equation}
\end{lemma}

\begin{proof}
If $\det B=0$, the choice $q=0$ gives $c\geq0$ and
$\adj(B)\psd$, so~\eqref{eq:adj-certificate} follows.  If $\det B>0$, then
$B\pd$ and the block condition is $q^*B^{-1}q\leq1$.  Given $y\neq0$, set
\[
 q=\frac{y}{\sqrt{y^*B^{-1}y}}.
\]
The hypothesis, multiplied by $y^*B^{-1}y$, says
$y^*(K+cB^{-1})y\geq0$.  Hence $K+cB^{-1}\psd$.  Multiplication by
$\det B>0$ and the identity $\adj(B)=(\det B)B^{-1}$ prove the claim.
\end{proof}

Applied to~\eqref{eq:c0}--\eqref{eq:K}, Lemma~\ref{lem:ellipsoid} gives the
positive-semidefinite certificate
\begin{equation}\label{eq:a4-certificate}
 (\det B(u,v,w))K(u,v,w)+c_0(u,v,w)\adj(B(u,v,w))\psd.
\end{equation}

\section{Completion of the proof}\label{sec:completion}

We now combine the geometric inequality with the finite character
classification.

\begin{proposition}[The two non-real $A_4$ characters]\label{prop:a4-characters}
For every $4\times4$ positive-semidefinite correlation matrix $A$,
\[
 d_{\chi_\omega}^{A_4}(A)\leq\per A,\qquad
 d_{\bar\chi_\omega}^{A_4}(A)\leq\per A.
\]
\end{proposition}

\begin{proof}
Write $A$ in the block form~\eqref{eq:completion} with
$B=B(a,b,d)$ and $q=(c,e,f)^T$.  Lemma~\ref{lem:a4-expansion} and
Theorem~\ref{thm:four-vector} give the first inequality.  An odd simultaneous
relabeling preserves positivity and the permanent and interchanges the two
non-real characters, giving the second.
\end{proof}

\begin{proof}[Proof of Theorem~\ref{thm:main}]
By Lemma~\ref{lem:correlation-reduction}, it is enough to consider
correlation matrices.  Lemma~\ref{lem:subgroup-classification} and the
transport Lemma~\ref{lem:transport} reduce an arbitrary subgroup and
irreducible representation to one of the listed character cases.
Proposition~\ref{prop:elementary-cases} proves thirty-five of them and
Proposition~\ref{prop:a4-characters} proves the remaining two.
Lemma~\ref{lem:reality} shows that the compared quantities are real.  Finally,
Lemma~\ref{lem:scaling} returns from correlation matrices to the full
positive-semidefinite cone.
\end{proof}

\appendix

\section{Character data}\label{app:character-data}

This appendix records the non-cyclic character data used in the finite
character reduction.  For $S_4$, use the conjugacy-class ordering
\[
 1,\quad(12),\quad(12)(34),\quad(123),\quad(1234).
\]
Its irreducible character table is
\begin{equation}\label{eq:s4-table}
\begin{array}{c|rrrrr}
 &1&(12)&(12)(34)&(123)&(1234)\\\hline
{[4]}    &1& 1& 1& 1& 1\\
{[1^4]}  &1&-1& 1& 1&-1\\
{[31]}   &3& 1&-1& 0&-1\\
{[211]}  &3&-1&-1& 0& 1\\
{[22]}   &2& 0& 2&-1& 0
\end{array}
\end{equation}

For $A_4$, let $C_+$ and $C_-$ be the two conjugacy classes of
three-cycles.  With $\omega$ as in~\eqref{eq:omega}, the table is
\begin{equation}\label{eq:a4-table}
\begin{array}{c|rrrr}
 &1&(12)(34)&C_+&C_-\\\hline
1&1&1&1&1\\
\chi_\omega&1&1&\omega&\omega^2\\
\bar\chi_\omega&1&1&\omega^2&\omega\\
\chi_3&3&-1&0&0.
\end{array}
\end{equation}

For $D_8=\langle r,s:r^4=s^2=1,srs=r^{-1}\rangle$, the four linear
characters take independent signs on $r$ and $s$, while its two-dimensional
character is $2,-2,0,0,0$ on
\[
 1,\quad r^2,\quad\{r,r^3\},\quad\{s,r^2s\},\quad\{rs,r^3s\}.
\]
The standard character tables of $C_m$, $C_2^2$, and $S_3$ supply the
remaining cases listed in Section~\ref{sec:elementary-cases}.

\section{Exact scalar and Gram certificates}\label{app:sos}

This appendix records the non-obvious polynomial certificates used in
Lemma~\ref{lem:hessian-certs}.  They contain no numerical approximation.
Recall $\nu=X^2+Y^2/3$.

\subsection{Direct square completions}

The coefficient $B_2$ has the exact decomposition
\begin{align*}
 B_2={}&38\left(\nu+\frac{39}{19}X-\frac{21}{19}Y+\frac32\right)^2
 +\frac{448}{19}\left(Y+\frac3{32}X-\frac{1083}{448}\right)^2\\
 &+\frac{3483}{16}\left(X+\frac{2137}{2322}\right)^2
 +\frac{131797}{5418}>0.
\end{align*}
Likewise
\begin{align*}
 B_3+252={}&6\left(\nu+\frac53X-2Y-\frac12\right)^2
 +24\left(Y-\frac{11}{12}X-\frac{15}{4}\right)^2\\
 &+\frac{139}{6}\left(X-\frac{249}{139}\right)^2
 +\frac{1851}{278}>0.
\end{align*}
Further elementary completions give
\[
 B_1\geq\frac{112}{3}\nu\geq0,\qquad
 B_4\geq\frac{132}{7}>0.
\]

\subsection{The conditional middle-coefficient certificate}

Set
\[
 \Sigma=4B_2B_4+252B_3.
\]
An exact Gram decomposition is
\begin{equation}\label{eq:sigma-gram}
 \Sigma=2128\nu\left(\nu+\frac{387}{133}X-
 \frac{261}{133}Y+\frac{77}{20}\right)^2
 +\frac1{59850}m^TQm,
\end{equation}
where $m=(1,Y,Y^2,X,XY,X^2)^T$ and
\begingroup
\scriptsize
\[
Q=\mat{
8739057600&-4231634400&478800000&6265576800&-1795500000&47880000\\
-4231634400&2519799114&-385003080&-2526627600&807975000&239400000\\
478800000&-385003080&79782560&281630160&-81252000&-35910000\\
6265576800&-2526627600&281630160&10931106942&-2359670040&2303554680\\
-1795500000&807975000&-81252000&-2359670040&685371360&-243756000\\
47880000&239400000&-35910000&2303554680&-243756000&1122611040
}.
\]
\endgroup
Exact rational $LDL^T$ elimination has positive pivots
\begin{align*}
&8739057600,\quad \frac{79557301866}{169},\quad
\frac{1112696849800960}{299087601},\\
&\frac{2624099206828303859949}{869294413907},\quad
\frac{2531308970961028469920080}{115380521779374043},\\
&\frac{1406694633431943161674826341030}
{10547120712337618624667}.
\end{align*}
Thus $Q\pd$ and $m^TQm>0$ because the first coordinate of $m$ is one.
Equation~\eqref{eq:sigma-gram} proves $\Sigma>0$.

If $B_3\geq0$, then $B_2+aB_3+a^2B_4>0$ is immediate.  If $B_3<0$, the
identity
\[
 4B_2B_4-B_3^2=(4B_2B_4+252B_3)+(-B_3)(B_3+252)
\]
shows that the discriminant of the quadratic
$B_2+aB_3+a^2B_4$ is negative.  Hence that quadratic is positive for every
$a>0$.  This completes the proof that $N_1>0$.

\section{Formal verification}

The theorem has been formalized in Lean 4.19.0 with Mathlib 4.19.0.  The
formal development verifies the subgroup and character enumeration, all
cycle expansions, the residual normalization identity, the rational Gram
certificate, and the final passage from correlation matrices to the full
positive-semidefinite cone.  The main statement is in the module
\[
 \texttt{NormalizedPermanentalDominance}.
\]
A detailed module map and an axiom audit are included with the source.  The
formalization is supplementary: the
identities and sign certificates needed for the mathematical proof have been
displayed above rather than delegated to the proof assistant.

The formal source used here was extracted from the immutable merge commit
\href{https://github.com/EndlessLethe/math/commit/c5e4d8c61f08d05ffe7c28bfa1ea732d8827b485}
{\texttt{c5e4d8c61f08}} and is mirrored in the companion repository.

\section*{Declaration of competing interest}
The author declares that there are no known competing financial interests or
personal relationships that could have appeared to influence the work
reported in this paper.

\section*{Data and code availability}
No external data were used in this study.  The accompanying Lean
formalization and the source of this manuscript are available from the
\href{https://github.com/EndlessLethe/PermanentalDominanceConjecture}
{companion repository}; the immutable formalization provenance is the commit
identified in the preceding section.

\bibliographystyle{elsarticle-num}
\bibliography{references}

\end{document}
